\documentclass{article}

\usepackage[utf8]{inputenc}
\usepackage[T5]{fontenc}
\usepackage[vietnamese]{babel}
\usepackage[final]{neurips_2024}
\usepackage{hyperref}
\usepackage{url}
\usepackage{booktabs}
\usepackage{amsfonts}
\usepackage{amsmath}
\usepackage{nicefrac}
\usepackage{microtype}
\usepackage{xcolor}
\usepackage{graphicx}
\usepackage{float}
\usepackage{subcaption}
\usepackage{enumitem}
\usepackage{array}
\usepackage{tabularx}
\usepackage{multirow}

\definecolor{propblue}{RGB}{0, 82, 155}
\hypersetup{colorlinks=true, linkcolor=propblue, urlcolor=propblue, citecolor=propblue}

\title{%
  \textbf{Báo cáo Thực nghiệm --- Yêu cầu 2 \& 3}\\[4pt]
  {\large Đề xuất và Đánh giá Mô hình EfficientNet-B0 + CBAM + ASL}\\
  {\large cho Bài toán Phân loại Ảnh Đa nhãn trên MS COCO}
}

\author{%
  Phan Huỳnh Châu Thịnh \quad MSSV: 23122019\\
  Hoàng Văn Sang \quad MSSV: 23120350\\
  \texttt{\{23122019, 23120350\}@student.hcmus.edu.vn}
}

\begin{document}
\maketitle

\begin{center}
  {\small\textit{Khoa Công nghệ Thông tin, Trường ĐHKH Tự nhiên, ĐHQG TP.HCM}}
\end{center}
\vspace{4pt}\noindent\rule{\linewidth}{0.8pt}

%% ===================================================================
\begin{abstract}
Báo cáo này trình bày \textbf{Yêu cầu 2} (đề xuất mô hình) và \textbf{Yêu cầu 3}
(thực nghiệm, phân tích, đánh giá) cho bài toán phân loại ảnh đa nhãn.
Chúng tôi đề xuất kiến trúc kết hợp \textbf{EfficientNet-B0} làm backbone,
\textbf{CBAM} làm neck tinh chỉnh đặc trưng, và \textbf{Asymmetric Loss (ASL)}
xử lý mất cân bằng nhãn (mô hình Exp~C).
Trên tập MS~COCO subset ($16{,}000$ ảnh train / $1{,}000$ val / $3{,}952$ test),
Exp~C đạt Val~mAP \textbf{0.6364} (best epoch~8), Test~mAP \textbf{0.6537}.
Overfitting gap (Train~mAP $0.9168$ $-$ Test~mAP $0.6537$ $= 0.263$) là thách thức
cần giải quyết, dẫn đến đề xuất cải tiến Exp~G (Yêu cầu~4).
\end{abstract}

%% ===================================================================
\section{Yêu cầu 2 --- Đề xuất Mô hình}

\subsection{Phân tích bài toán và động cơ đề xuất}

Bài toán phân loại ảnh đa nhãn yêu cầu mô hình dự đoán đồng thời nhiều nhãn
cho một ảnh đầu vào. Trên MS~COCO (80 lớp), mỗi ảnh trung bình chỉ chứa
$2$--$3$ nhãn dương tính, tức tỉ lệ positive/negative cực thấp ($\approx 1{:}37$).
Hai vấn đề cốt lõi cần giải quyết:

\begin{enumerate}[leftmargin=2em]
  \item \textbf{Mất cân bằng nhãn nghiêm trọng:} Binary Cross-Entropy (BCE) truyền
        thống không phân biệt gradient từ easy negatives và hard positives,
        dẫn đến mô hình thiên lệch về phía âm tính.
  \item \textbf{Đặc trưng không gian chưa được khai thác:} Trong ảnh đa nhãn,
        mỗi nhãn thường gắn với một vùng không gian cụ thể; backbone thông thường
        không có cơ chế định vị vùng này trước khi tổng hợp đặc trưng.
\end{enumerate}

\subsection{Kiến trúc đề xuất (Exp~C)}

Chúng tôi đề xuất kiến trúc feed-forward gồm 5 giai đoạn:

\begin{equation}
  \text{Ảnh} \;\xrightarrow{\text{EfficientNet-B0}}\; F
  \;\xrightarrow{\text{CBAM}}\; \tilde{F}
  \;\xrightarrow{\text{GAP}}\; \mathbf{v}
  \;\xrightarrow{\text{FC}+\text{Dropout}}\; \mathbf{z}
  \;\xrightarrow{\text{ASL}}\; \mathcal{L}
\end{equation}

\noindent Trong đó $F \in \mathbb{R}^{B \times 1280 \times 7 \times 7}$
là feature map từ backbone, $\tilde{F}$ là feature map sau CBAM,
$\mathbf{v} \in \mathbb{R}^{1280}$ là vector đặc trưng sau GAP,
$\mathbf{z} \in \mathbb{R}^{80}$ là vector logit trước sigmoid.

\paragraph{Backbone --- EfficientNet-B0.}
EfficientNet-B0~\cite{tan2019efficientnet} sử dụng compound scaling để cân bằng
đồng thời chiều sâu, chiều rộng, và độ phân giải. Với chỉ $3.63\text{M}$ tham số
(so với $23.67\text{M}$ của ResNet50), EfficientNet-B0 nhẹ hơn $6.5\times$,
phù hợp với giới hạn phần cứng Kaggle~T4 (16\,GB VRAM).

\paragraph{Neck --- CBAM.}
CBAM (Convolutional Block Attention Module)~\cite{woo2018cbam} áp dụng
tuần tự hai cơ chế chú ý:

\begin{itemize}[leftmargin=1.8em]
  \item \textbf{Channel Attention} học trọng số $\mathbf{s}_c \in (0,1)^C$
        để nhấn mạnh kênh đặc trưng nào quan trọng (``what''):
        \[
          \mathbf{s}_c = \sigma\!\bigl(\text{MLP}(\text{AvgPool}(F))
          + \text{MLP}(\text{MaxPool}(F))\bigr)
        \]
  \item \textbf{Spatial Attention} tạo bản đồ chú ý $\mathbf{M}_s \in (0,1)^{H\times W}$
        để định vị vùng nào trong ảnh đáng chú ý (``where''):
        \[
          \mathbf{M}_s = \sigma\!\bigl(\text{Conv}_{7\times7}(
          [\text{AvgPool}_c(F'),\, \text{MaxPool}_c(F')])\bigr)
        \]
\end{itemize}

\paragraph{Loss --- Asymmetric Loss (ASL).}
ASL~\cite{ridnik2021asl} xử lý riêng nhãn dương ($y=1$) và âm ($y=0$)
với hệ số focusing khác nhau $\gamma_+ < \gamma_-$:
\begin{equation}
  L_{\text{ASL}} = \begin{cases}
    (1-p)^{\gamma_+}\log p & y=1 \\
    (p_{\text{shift}})^{\gamma_-}\log(1-p_{\text{shift}}) & y=0
  \end{cases}
\end{equation}
Với $p_{\text{shift}} = \max(p-m, 0)$, $m=0.05$, $\gamma_+=0$, $\gamma_-=4$.
Cơ chế này triệt tiêu hoàn toàn gradient của easy negatives ($p < m$),
giải quyết triệt để vấn đề mất cân bằng nhãn.

\subsection{Bản chất toán học: Mô hình là chuỗi các phép chiếu W}

\textbf{Nhìn từ góc độ đại số tuyến tính,} toàn bộ pipeline có thể hiểu là
chuỗi các phép biến đổi vector đặc trưng:

\begin{enumerate}[leftmargin=2em]
  \item \textbf{Backbone} học các ma trận tích chập $\mathbf{W}_\ell^{\text{conv}}$
        tại mỗi layer $\ell$, biến đổi ảnh thô thành feature map $F$ ở không gian
        đặc trưng ngữ nghĩa cao.
  \item \textbf{CBAM Channel} học ma trận MLP $\mathbf{W}_{\text{ch}}$ chiếu
        vector kênh về không gian trọng số để chọn lọc kênh quan trọng.
  \item \textbf{CBAM Spatial} học kernel conv $\mathbf{W}_{\text{sp}} \in \mathbb{R}^{1\times2\times7\times7}$
        để gắn trọng số không gian, định vị đối tượng.
  \item \textbf{FC Head} học ma trận $\mathbf{W}_{\text{fc}} \in \mathbb{R}^{80 \times 1280}$
        chiếu vector đặc trưng $\mathbf{v}$ về không gian nhãn: logit thứ $i$ =
        $\mathbf{w}_i^\top \mathbf{v}$ đo mức độ tương đồng giữa $\mathbf{v}$
        và prototype của lớp $i$.
\end{enumerate}

\noindent Sự phân biệt giữa các lớp được đảm nhiệm bởi \textbf{ma trận FC}:
hàng $i$ của $\mathbf{W}_{\text{fc}}$ là ``prototype vector'' của lớp $i$ trong
không gian đặc trưng. Mô hình phân loại đúng khi các prototype đủ xa nhau,
được tối ưu bởi ASL thông qua gradient amplification trên hard samples.

%% ===================================================================
\section{Yêu cầu 3 --- Thực nghiệm, Phân tích và Đánh giá}

\subsection{Thiết lập thực nghiệm}

\begin{table}[H]
  \caption{Các kịch bản ablation study (6 experiments).}
  \label{tab:ablation_setup}
  \centering
  \begin{tabular}{clccc}
    \toprule
    \textbf{Exp} & \textbf{Mô tả} & \textbf{Backbone} & \textbf{CBAM} & \textbf{Loss} \\
    \midrule
    A & Baseline              & ResNet50        & Không & BCE   \\
    B & ASL thay BCE          & ResNet50        & Không & ASL   \\
    C & \textbf{Mô hình đề xuất} & EfficientNet-B0 & Có & ASL   \\
    D & Focal Loss            & ResNet50        & Không & Focal \\
    E & ResNet + CBAM + ASL   & ResNet50        & Có    & ASL   \\
    F & EfficientNet + ASL    & EfficientNet-B0 & Không & ASL   \\
    \bottomrule
  \end{tabular}
\end{table}

\noindent \textbf{Dữ liệu:} MS~COCO 2017 subset stratified:
$16{,}000$ ảnh train / $1{,}000$ val / $3{,}952$ test.\\
\textbf{Tiền xử lý:} resize $224\times224$, RandomHorizontalFlip,
ColorJitter, chuẩn hóa ImageNet ($\mu$=\{0.485,0.456,0.406\},
$\sigma$=\{0.229,0.224,0.225\}).\\
\textbf{Môi trường:} Kaggle Notebook, GPU NVIDIA Tesla T4 (16\,GB VRAM).\\
\textbf{Siêu tham số:} AdamW, lr=$3\times10^{-4}$, OneCycleLR,
batch=64, epochs=20, dropout=0.3.

\subsection{Kết quả định lượng (số liệu thực nghiệm)}

\begin{table}[H]
  \caption{Kết quả toàn diện: Val~mAP (best epoch), Train~mAP, Test~mAP,
           Overfitting Gap (Train $-$ Test), và các chỉ số F1.}
  \label{tab:full_results}
  \centering
  \small
  \begin{tabular}{lccccccc}
    \toprule
    \textbf{Exp} & \textbf{Val mAP} & \textbf{Train mAP}
      & \textbf{Test mAP} & \textbf{Gap} & \textbf{Val F1} & \textbf{Test F1} \\
    \midrule
    A: ResNet50+BCE      & 0.6908 & 0.8706 & 0.7081 & 0.1625 & 0.6201 & 0.6486 \\
    B: ResNet50+ASL      & 0.7107 & 0.8739 & 0.7228 & 0.1511 & 0.5886 & 0.6159 \\
    D: ResNet50+Focal    & 0.7009 & 0.8665 & 0.7100 & 0.1565 & 0.6286 & 0.6467 \\
    E: ResNet50+CBAM+ASL & 0.7129 & 0.9310 & 0.7223 & 0.2087 & 0.6167 & 0.6274 \\
    F: EffNet+ASL        & 0.6326 & 0.8132 & 0.6463 & 0.1669 & 0.5460 & 0.5598 \\
    \textbf{C: EffNet+CBAM+ASL} & \textbf{0.6364} & \textbf{0.9168} & \textbf{0.6537} & \textbf{0.2631} & \textbf{0.5651} & \textbf{0.5768} \\
    \bottomrule
  \end{tabular}
\end{table}

\noindent \textit{Ghi chú:} Train~mAP được tính bằng cách chạy evaluate trên
toàn bộ tập train ($16{,}000$ ảnh) với checkpoint best.pth.
Val~mAP được ghi lại từ log.json (best epoch trong quá trình training).

\subsection{Phân tích: Tại sao không đạt 100\%?}

\subsubsection{Mất cân bằng nhãn nghiêm trọng}

Trên MS~COCO subset ($16{,}000$ ảnh, 80 lớp):
\begin{itemize}[leftmargin=1.8em]
  \item Lớp phổ biến nhất (\textit{person}): $\approx 7{,}000$ positive samples.
  \item Lớp hiếm nhất (\textit{hair drier}): $< 150$ positive samples.
  \item Tỉ lệ positive trung bình: $\sim 2.7\%$ mỗi lớp ($\approx 37{:}1$).
\end{itemize}
Dù ASL xử lý tốt hơn BCE, các lớp hiếm gặp vẫn đạt AP rất thấp
(ví dụ: \textit{hair drier} AP $< 0.12$).

\subsubsection{Giới hạn dữ liệu và backbone}

Chỉ dùng $16{,}000$ ảnh (thay vì $118{,}000$ ảnh toàn bộ COCO), độ bao phủ
phân phối thực tế không đầy đủ. EfficientNet-B0 với $3.63\text{M}$ tham số
có thể bị \textit{under-parameterized} cho tác vụ 80 lớp phức tạp
(so với ResNet50 với $23.67\text{M}$ tham số đạt Test~mAP~$0.7228$).

\subsubsection{Nhập nhằng giữa các lớp (Class Confusion)}

Nhiều nhãn đồng xuất hiện thường xuyên (\textit{person + bicycle, chair + dining table}),
khiến mô hình khó phân biệt boundary quyết định.
Threshold cố định $\theta = 0.5$ không tối ưu cho mọi lớp ---
đây là lý do chính dẫn đến Macro-F1 thấp trong Exp~C.

\subsection{Train cao hơn Test: Hiện tượng Overfitting}

\begin{table}[H]
  \caption{Phân tích overfitting gap.}
  \centering
  \small
  \begin{tabular}{lcc}
    \toprule
    \textbf{Exp} & \textbf{Gap (Train$-$Test)} & \textbf{Nhận xét} \\
    \midrule
    B: ResNet50+ASL      & 0.1511 & Overfitting nhẹ \\
    A: ResNet50+BCE      & 0.1625 & Overfitting nhẹ \\
    D: ResNet50+Focal    & 0.1565 & Overfitting nhẹ \\
    F: EffNet+ASL        & 0.1669 & Overfitting nhẹ \\
    E: ResNet50+CBAM+ASL & 0.2087 & Overfitting trung bình \\
    C: EffNet+CBAM+ASL   & 0.2631 & \textbf{Overfitting nặng} \\
    \bottomrule
  \end{tabular}
\end{table}

\textbf{Lý giải tại sao Train~mAP $>$ Test~mAP:}
\begin{enumerate}[leftmargin=2em]
  \item \textbf{Distribution shift:} Mô hình ghi nhớ đặc trưng cụ thể của
        $16{,}000$ ảnh train, không generalize hoàn toàn sang tập test.
  \item \textbf{CBAM ``học nhiễu'':} Với dataset nhỏ, CBAM có thể học các
        attention map đặc thù với train set (texture/color bias) thay vì
        học đặc trưng tổng quát --- đây là lý do Exp~C có gap lớn nhất
        ($0.2631$) dù backbone nhẹ hơn.
  \item \textbf{Capacity mismatch:} Exp~E (ResNet50+CBAM) có gap $0.2087$,
        cho thấy CBAM kết hợp backbone lớn cũng dễ overfit hơn khi không
        có thêm regularization.
  \item \textbf{Threshold \& metric:} Train~mAP được tính với threshold tối ưu
        tìm trên train set, trong khi Test dùng threshold cố định $0.5$.
\end{enumerate}

\subsection{Phân tích per-class AP --- Lớp nào khó nhất?}

\begin{table}[H]
  \caption{10 lớp có AP thấp nhất (Exp~C, Val set) và nguyên nhân.}
  \label{tab:perclass}
  \centering
  \small
  \begin{tabular}{lcp{5cm}}
    \toprule
    \textbf{Lớp} & \textbf{AP} & \textbf{Nguyên nhân} \\
    \midrule
    hair drier   & $<0.12$ & Rất hiếm ($<150$ mẫu), thường bị che khuất \\
    toothbrush   & $<0.15$ & Nhỏ trong ảnh, nhiều màu sắc khác nhau \\
    parking meter& $<0.20$ & Nền đô thị gây nhiễu, hình dạng dài-mảnh \\
    scissors     & $<0.20$ & Nhỏ, nhiều tư thế, co-occurs với bàn \\
    snowboard    & $<0.30$ & Co-occurs với \textit{person}, bối cảnh tuyết trắng \\
    cell phone   & $<0.40$ & Bị che tay, góc chụp đa dạng \\
    mouse        & $<0.45$ & Nhỏ, phụ thuộc context bàn làm việc \\
    skis         & $<0.45$ & Co-occurs với \textit{snowboard}, tuyết trắng \\
    remote       & $<0.47$ & Nhỏ, hình dạng đơn giản, dễ nhầm \\
    tie          & $<0.47$ & Bị che bởi quần áo, phụ thuộc \textit{person} \\
    \bottomrule
  \end{tabular}
\end{table}

\noindent \textbf{Nhận xét quan trọng:} Spatial Attention ($7\times7$)
của CBAM không đủ độ phân giải để định vị vật thể rất nhỏ
(hair drier, toothbrush, scissors). Đây là hạn chế cấu trúc cần cải tiến.

\subsection{Phân tích 10 mẫu test tiêu biểu --- Trường hợp sai}

Phân tích định tính trên các mẫu test có lỗi cao nhất (FP+FN nhiều nhất):

\begin{table}[H]
  \caption{Phân loại các trường hợp dự đoán sai điển hình của Exp~C.}
  \label{tab:error}
  \centering
  \small
  \begin{tabular}{p{2.5cm}p{3cm}p{4cm}p{3.5cm}}
    \toprule
    \textbf{Loại lỗi} & \textbf{Ví dụ} & \textbf{Nguyên nhân} & \textbf{Hướng cải tiến}\\
    \midrule
    False Negative (FN) &
      \textit{hair drier} trong ảnh nhưng bỏ sót &
      $<150$ mẫu train, vật nhỏ, bị che &
      Oversample, focal loss mạnh hơn \\
    False Positive (FP) &
      Predict \textit{tie} khi không có &
      Co-occurrence cao với \textit{person} trong train &
      Debiasing, negative sampling \\
    Nhầm lớp tương đồng &
      Nhầm \textit{skis} $\leftrightarrow$ \textit{snowboard} &
      Hai lớp cùng bối cảnh, hình dạng gần nhau &
      Contrastive loss \\
    Mẫu ngoài phân phối &
      Ảnh ban đêm, góc cực đoan &
      Train chủ yếu trên ảnh ban ngày &
      Data augmentation đa dạng hơn \\
    \midrule
    \multicolumn{4}{l}{\textit{Nhận xét:} Threshold $0.5$ cố định gây FP cho lớp
    co-occurs cao (\textit{person+tie})} \\
    \multicolumn{4}{l}{\textit{và FN cho lớp hiếm (\textit{hair drier}).
    $\Rightarrow$ Động cơ cho Adaptive Threshold (Exp~G).}} \\
    \bottomrule
  \end{tabular}
\end{table}

\subsection{Phân tích đường cong training}

\paragraph{Exp~A (ResNet50+BCE):}
Train loss giảm đều từ $\approx0.114$ (epoch~2) xuống $0.034$ (epoch~20, hội tụ).
Val mAP tăng từ $0.422$ lên $0.691$. Hội tụ rõ từ epoch~10.

\paragraph{Exp~C (EfficientNet+CBAM+ASL):}
Train loss khởi đầu cao ($\approx 1.12$, epoch~2) do ASL scale khác BCE,
giảm dần xuống $0.943$ (epoch~8 --- best val mAP). Val mAP tăng từ $0.557$
(epoch~2) lên $0.636$ (epoch~8), sau đó dao động và không tăng thêm nhiều.
\textbf{Nguyên nhân plateau:} EfficientNet-B0 với $3.63\text{M}$ tham số
đã học hết thông tin có thể trích xuất từ $16{,}000$ ảnh train.

\paragraph{Exp~B (ResNet50+ASL) --- Baseline tốt nhất:}
Đạt Val~mAP $0.7107$ (epoch~8), Train~mAP $0.8739$, Test~mAP $0.7228$.
ResNet50 với $23.67\text{M}$ tham số có đủ capacity cho tác vụ 80 lớp,
và ASL giải quyết tốt mất cân bằng nhãn mà không cần CBAM.

\subsection{Nhận xét tổng kết (Yêu cầu 3)}

\begin{enumerate}[leftmargin=2em]
  \item \textbf{Về mAP:} Exp~B (ResNet50+ASL) dẫn đầu với Test~mAP $0.7228$.
        Exp~C (mô hình đề xuất) đạt $0.6537$ --- thấp hơn Exp~B do
        EfficientNet-B0 bị under-parameterized cho 80 lớp.
  \item \textbf{Về F1:} Exp~A (ResNet50+BCE) và Exp~D (ResNet50+Focal)
        đạt Test~F1 cao nhất ($0.6486$ và $0.6467$) nhờ threshold~$0.5$
        phù hợp với BCE/Focal. Exp~C có F1 thấp ($0.5768$) --- đây là
        vấn đề cần cải tiến trong Exp~G.
  \item \textbf{Về overfitting:} Exp~C có gap lớn nhất ($0.2631$),
        cho thấy CBAM trên dataset nhỏ dễ học noise hơn là đặc trưng tổng quát.
  \item \textbf{Kết luận:} Kết hợp CBAM+EfficientNet không cải thiện mAP
        so với baseline ResNet50+ASL, nhưng tạo ra \textit{cấu trúc nhẹ hơn}
        ($3.63\text{M}$ vs $23.67\text{M}$ tham số) với potential lớn nếu
        giải quyết được vấn đề threshold (dẫn đến Exp~G).
\end{enumerate}

%% ===================================================================
\bibliographystyle{plain}
\bibliography{references}

\end{document}
